CKKS 方案支持复数上的 SIMD 近似算术，实数是其中的特例，适用于机器学习等数值计算任务。BFV、BGV 和 TFHE 的原生明文空间是精确的有限集合，但通过适当的编码或电路也能实现定点或浮点运算；它们与 CKKS 的区别在于原生算术及误差的含义。~\cite{ckks-original}

对于二次幂 $n\geq 2$，应区分实代数 $\mathcal{R}_{\langle n\rangle}=\mathbb{R}[X]/(X^n+1)$、整数环 $R=\mathbb{Z}[X]/(X^n+1)$ 和密文环 $R_q=R/qR$。典范嵌入的逆映射将 $n/2$ 个复数槽位映为多项式 $M\in\mathcal{R}_{\langle n\rangle}$。随后，编码将其系数乘以 $\Delta>0$ 并舍入，得到 $R$ 中的整数多项式，再在模 $q$ 下加密。


$ $

\noindent 总体而言，CKKS 的加密过程如下：

\begin{enumerate}
\item \underline{\textsf{编码\textsubscript{1}}：}将目标输入复数编码为实数
\item \underline{\textsf{编码\textsubscript{2}}：}将实数编码为整数
\item \underline{\textsf{加密}：}使用 RLWE 加密整数
\end{enumerate}
得到的 RLWE 密文支持同态加法和乘法。

$ $

\noindent 所有同态运算完成后，CKKS 的解密过程如下：

\begin{enumerate}
\item \underline{\textsf{解密}：}将 RLWE 密文解密为明文整数
\item \underline{\textsf{解码\textsubscript{1}}：}将整数解码为实数
\item \underline{\textsf{解码\textsubscript{2}}：}将实数解码为复数
\end{enumerate}

$ $

CKKS 将放大后的逆嵌入舍入到整数格时，通常会产生编码误差。加密噪声以及后续算术、密钥切换、重缩放和函数近似产生的误差也会影响解码结果。这并非专门由平方根产生的误差，某些输入也能精确编码。与 BFV 的精确明文解密不同，一般的 CKKS 解密会保留数值误差。

CKKS 和 BFV 共用环算术与密钥切换技术，但编码、解码、尺度管理、乘法和自举必须遵循各自不同的明文含义。两者都使用环：$\mathbb{C}$ 本身是域，因此也是环。CKKS 的关键区别是特征为零的近似算术，而不是在明文模数下进行精确算术。



\begin{tcolorbox}[
    title = \textbf{所需背景知识},    % box title
    colback = white,    % light background; tweak to taste
    colframe = black,  % frame colour
    boxrule = 0.8pt,     % line thickness
    left = 1mm, right = 1mm, top = 1mm, bottom = 1mm % inner padding
]

\begin{itemize}
\item \autoref{sec:modulo}: \nameref{sec:modulo}
\item \autoref{sec:group}: \nameref{sec:group}
\item \autoref{sec:field}: \nameref{sec:field}
\item \autoref{sec:order}: \nameref{sec:order}
\item \autoref{sec:polynomial-ring}: \nameref{sec:polynomial-ring}
\item \autoref{sec:decomp}: \nameref{sec:decomp}
\item \autoref{sec:roots}: \nameref{sec:roots}
\item \autoref{sec:cyclotomic}: \nameref{sec:cyclotomic}
\item \autoref{sec:matrix}: \nameref{sec:matrix}
\item \autoref{sec:euler}: \nameref{sec:euler}
\item \autoref{sec:modulus-rescaling}: \nameref{sec:modulus-rescaling}
\item \autoref{sec:chinese-remainder}: \nameref{sec:chinese-remainder}
\item \autoref{sec:taylor-series}: \nameref{sec:taylor-series}
\item \autoref{sec:polynomial-interpolation}: \nameref{sec:polynomial-interpolation}
\item \autoref{sec:ntt}: \nameref{sec:ntt}
\item \autoref{sec:lattice}: \nameref{sec:lattice}
\item \autoref{sec:rlwe}: \nameref{sec:rlwe}
\item \autoref{sec:glwe}: \nameref{sec:glwe}
\item \autoref{sec:glev}: \nameref{sec:glev}
\item \autoref{sec:glwe-add-cipher}: \nameref{sec:glwe-add-cipher}
\item \autoref{sec:glwe-add-plain}: \nameref{sec:glwe-add-plain}
\item \autoref{sec:glwe-mult-plain}: \nameref{sec:glwe-mult-plain}
\item \autoref{subsec:modulus-switch-rlwe}: \nameref{subsec:modulus-switch-rlwe}
\item \autoref{sec:glwe-key-switching}: \nameref{sec:glwe-key-switching}
\item \autoref{sec:bfv}: \nameref{sec:bfv}
\end{itemize}
\end{tcolorbox}

\clearpage


\subsection{编码与解码}
\label{subsec:ckks-encoding-decoding}


CKKS 的编码与解码在本质上非常接近 BFV 的批量编码方案。BFV 基于更新后的 $\hathat W$ 和 $\hathat W^*$ 矩阵（\autoref{subsec:bfv-rotation} 中的总结~\ref*{subsubsec:bfv-rotation-summary}）设计批量编码方案（\autoref{subsec:bfv-enc-dec} 中的总结~\ref*{subsec:bfv-enc-dec}）。具体而言，BFV 在 $X^n+1$ 的各个根，即本原 $(\mu=2n)$ 次单位根处对多项式求值，从而将多项式解码为输入槽位向量（即 $\vec{v} = \hathat W^* \cdot \vec{m}$）；反过来，通过这一运算的逆运算将输入槽位向量编码为多项式（即 $\vec{m} = n^{-1} \cdot \hathat W \cdot I_n^R \cdot \vec{v}$）。这一编码与解码方案以总结~\ref*{subsec:poly-vector-transformation-complex}（\autoref{subsec:poly-vector-transformation-complex}）为基础，后者在环（有限域）上的 $n$ 槽位向量与 $(n-1)$ 次（或更低次）多项式之间构造了如下同构映射：

$\sigma: f(x) \in \mathbb{Z}_t[X] / F(X) \text{ } \longrightarrow \text{ } (f(\omega^1), f(\omega^3)), \cdots, f(\omega^{2n-1})) \in \mathbb{Z}_t^n$

其中，$\omega = g^{\frac{t - 1}{2n}}$ 是定义在模素数 $t$ 的环上的第 $(\mu=2n)$ 个分圆多项式 $X^n + 1$ 的根，即本原 $(\mu=2n)$ 次单位根。

$ $

CKKS 在复单位根上使用相同的求值与插值原理。复数也构成环；这里的完整求值向量属于具有共轭对称性的实代数 $\mathbb{\hat C}^n$。总结~\ref*{subsec:poly-vector-transformation-complex}（\autoref{subsec:poly-vector-transformation-complex}）给出了如下典范嵌入：

$\sigma_c: f(X) \in \mathbb{R}[X]/(X^n + 1) \longrightarrow (f(\omega),f(\omega^3),f(\omega^5), \cdots, f(\omega^{2n-1})) \in \mathbb{\hat{C}}^{n} \text{ } (\longrightarrow \mathbb{C}^{\frac{n}{2}})$


其中，$\omega = e^{i\pi/n}$ 是定义在复数域上的第 $(\mu=2n)$ 个分圆多项式 $X^n + 1$ 的根，即本原 $(\mu=2n)$ 次单位根；$\mathbb{\hat{C}}^{n}$ 是一个特殊的 $n$ 维复向量空间，其后半部分元素是前半部分元素按逆序排列的共轭。$\mathbb{\hat{C}}^{n}$ 与 $\mathbb{{C}}^{\frac{n}{2}}$ 同构，因为 $\mathbb{\hat{C}}^{n}$ 的后半部分元素完全由前半部分元素确定。因此，映射 $\sigma_c$ 本质上是在 $\dfrac{n}{2}$ 槽位复向量 $\vec{v} \in \mathbb{C}^{\frac{n}{2}}$ 与 $(n-1)$ 次（或更低次）实系数多项式 $\mathbb{R}[X] / (X^n + 1)$ 之间的同构。所以，CKKS 的批量编码方案将一个 $\dfrac{n}{2}$ 槽位的复数输入向量编码为 $(n-1)$ 次（或更低次）的实系数多项式，而解码过程则与此相反。

此外，回顾 BFV 中的做法：为了支持输入向量槽位的同态旋转，我们将 $W$ 和 $W^T$ 更新为 $\hathat W$ 和 $\hathat W^*$（见 \autoref{subsubsec:bfv-rotation-summary} 中的总结~\ref*{subsubsec:bfv-rotation-summary}）。同样，CKKS 批量编码方案使用 $\hathat W$ 和 $\hathat W^*$ 替代 $W$ 和 $W^T$，以支持同态旋转。因此，CKKS 批量编码方案中的同构映射更新为：


$\sigma_c: f(X) \in \mathbb{R}[X]/(X^n + 1) \longrightarrow \bm ( f(\omega^{J(0)}),f(\omega^{J(1)}),f(\omega^{J(2)}), \cdots, f(\omega^{J(\frac{n}{2} - 1)}), $

\textcolor{white}{$ f(X) \in \mathbb{R}[X]/(X^n + 1) \longrightarrow \bm ( $} $\cdots, f(\omega^{J_*(0)}),f(\omega^{J_*(1)}),f(\omega^{J_*(2)}), \cdots, f(\omega^{J_*(\frac{n}{2} - 1)}) \bm) \in \mathbb{\hat{C}}^{n} \text{ } (\longrightarrow \mathbb{C}^{\frac{n}{2}})$

其中，$J(h) = 5^h \bmod 2n$ 是用于辅助旋转的公式。


$ $

BFV 与 CKKS 的编码方案具有以下区别：

\begin{itemize}
\item \textbf{输入槽位值的类型：}BFV 的输入槽位值是 ${n}$ 个模 $t$ 的整数，它们被编码为 $n$ 个多项式系数（即模 $t$ 的整数）。而 CKKS 的输入槽位值是 $\dfrac{n}{2}$ 个复数，它们被编码为 $n$ 个多项式系数（即实数）。

\item \textbf{多项式系数的类型：} BFV 编码到 $\mathbb{Z}_t[X]/(X^n+1)$。CKKS 先插值得到实数系数，再将放大后的系数舍入为整数。两种方案实际加密的多项式都具有整数系数。

\item \textbf{尺度因子：} BFV 在密文模数 $q$ 下使用 $\Delta=\lfloor q/t\rfloor$。CKKS 根据目标精度、消息界和模数链选择 $\Delta$。正确解释密文要求包含全部累积误差的解密相位满足中心化不回绕界；系数可以为零或负数。在每一层都要求 $\|\Delta M+E\|_\infty<q_l/2$。

\item \textbf{编码精度：}对于 BFV，在解码过程中，将多项式系数缩小为 $\dfrac{ \Delta\cdot \vec{m} }{\Delta}$ 后能够保留输入值的精度
%as far as $\Delta$ is sufficiently large enough to round off the scaling error (as explained in \autoref{subsubsec:bfv-enc-dec-decoding1})
。而对于 CKKS，如果原始输入值的小数位数过多，尺度因子无法将所有小数位都左移到整数部分，那么缩小后的多项式系数就可能丢失精度。这意味着输入值的一些低位小数会被舍入，导致原始输入的精度损失。例如，设多项式系数为 $m_i = \dfrac{1}{3} = 0.33333\cdots$，尺度因子为 $\Delta = 100$。则缩放后的系数为 $\lceil \Delta m_i \rfloor = 33$，再缩小后得到 $\dfrac{33}{100} = 0.33$。由于 $0.33 \neq 0.33333\cdots$，CKKS 的编码与解码过程并不总能保证完全精确。正是由于这种编码误差，CKKS 被称为\textit{近似}加密方案。编码误差的影响可能随同态运算而增大，使误差幅度增加，解码结果逐渐偏离预期的精确值。减小 CKKS 编码误差的一种方式是增大 $\Delta$，使更多小数位左移到缩放后的整数部分。

\end{itemize}

$ $


\para{\boldmath$\vec{v}_{'} \in \mathbb{\hat C}^n$ 的结构：}注意，总结~\ref*{subsec:poly-vector-transformation-complex}（\autoref{subsec:poly-vector-transformation-complex}）中给出的 $\vec{v}_{'}$ 的原始解码方案为：

$\vec{v}_{'} = \bm{(} \text{ } M(\omega), \text{ } M(\omega^3), \text{ } M(\omega^5), \cdots, M(\omega^{2n-3}), \text{ } M(\omega^{2n-1}) \bm{)}$

它解码得到一个厄米向量：

$\vec{v}_{'} = (v_0, v_1, \cdots, v_{\frac{n}{2} - 1}, \overline v_{\frac{n}{2} - 1}, \cdots, \overline v_1, \overline v_0 )$

其后半部分元素是前半部分元素按逆序排列的共轭。

$ $

不过，将 $W$ 和 $W^T$ 替换为 $\hathat W$ 和 $\hathat W^*$ 后，上述解码方案改为以下支持同态旋转的形式：


$\vec{v}_{'} =  \bm{(} \text{ }
M(\omega^{J(0)}), \text{ } M(\omega^{J(1)}), \text{ } M(\omega^{J(2)}), \cdots,  M(\omega^{J(\frac{n}{2}-1)}), \text{ } M(\omega^{J_*(0)}), \text{ } M(\omega^{J_*(1)}), \cdots,  M(\omega^{J_*(\frac{n}{2}-1)}) \text{ } \bm{)}$

$\textcolor{white}{\vec{v}_{'}} =  \bm{(} \text{ }
M(\omega^{J(0)}), \text{ } M(\omega^{J(1)}), \text{ } M(\omega^{J(2)}), \cdots,  M(\omega^{J(\frac{n}{2}-1)}), \text{ } M(\overline\omega^{J(0)}), \text{ } M(\overline\omega^{J(1)}), \cdots,  M(\overline\omega^{J(\frac{n}{2}-1)}) \text{ } \bm{)}$

\textcolor{red}{ $\rhd$ 因为 $\omega^{-1} = (e^{\frac{i\pi}{n}})^{-1} = e^{\frac{-i\pi}{n}} = \overline\omega$，且给定 $J(h) = 5^h \bmod 2n$ 和 $J_*(h) = -5^h \bmod 2n$}

它解码得到如下\textit{正序}（而非逆序）厄米向量：

$\vec{v}_{'} = (v_0, v_1, \cdots, v_{\frac{n}{2} - 1}, \overline v_0, \overline v_1, \cdots, \overline v_{\frac{n}{2} - 1})$

其后半部分元素是前半部分元素按相同顺序排列的共轭。执行同态旋转时（将在 \autoref{subsec:ckks-rotation} 中说明），与 BFV 的同态旋转一样，$\vec{v}_{'}$ 的前半部分和后半部分元素分别在各自的组内循环旋转。



下面总结 CKKS 的编码与解码过程。它与 BFV 的编码与解码过程相似，后者见 \autoref{subsubsec:bfv-encoding-summary} 中的总结~\ref*{subsubsec:bfv-encoding-summary}：

\begin{tcolorbox}[title={\textbf{\tboxlabel{\ref*{subsec:ckks-encoding-decoding}} CKKS 的编码与解码}}]

\textbf{\underline{输入}：}一个 $\dfrac{n}{2}$ 槽位的复向量 $\vec{v} = (v_0, v_1, \cdots, v_{\frac{n}{2}-1}) \in \mathbb{C}^{\frac{n}{2}}$

\par\noindent\rule{\textwidth}{0.4pt}

\textbf{\underline{编码}：}

$ $

\begin{enumerate}
\item 将 $\vec{v}$ 转换（即同构变换）为如下 $n$ 槽位的\textit{正序}厄米向量 $\vec{v}_{'}$：

$\vec{v}_{'} = (v_0, v_1, \cdots, v_{\frac{n}{2}-1}, \overline v_0, \overline v_1, \cdots, \overline v_{\frac{n}{2}-1}) \in \mathbb{\hat C}^{n}$

\item 应用变换 $\vec{m} = \dfrac{\hathat W \cdot I_n^R \cdot \vec{v}_{'}}{n}$，将 $\vec{v}_{'}$ 转换为实数向量 $\vec{m}$

其中，$\hathat{W}$ 是按如下方式构造的 $n$ 维向量空间的一组基：

$ $

\noindent{\footnotesize{\noindent $\noindent\hathat W = \begin{bmatrix}
1 & 1 & \cdots & 1 & 1 & 1 & \cdots & 1\\
(\omega^{J(\frac{n}{2} - 1)}) & (\omega^{J(\frac{n}{2} - 2)}) & \cdots & (\omega^{J(0)}) & (\omega^{J_*(\frac{n}{2} - 1)}) & (\omega^{J_*(\frac{n}{2} - 2)}) & \cdots & (\omega^{J_*(0)})\\
(\omega^{J(\frac{n}{2} - 1)})^2 & (\omega^{J(\frac{n}{2} - 2)})^2 & \cdots & (\omega^{J(0)})^2 & (\omega^{J_*(\frac{n}{2} - 1)})^2 & (\omega^{J_*(\frac{n}{2} - 2)})^2 & \cdots & (\omega^{J_*(0)})^2 \\
\vdots & \vdots & \ddots & \vdots & \vdots & \ddots & \vdots & \vdots \\
(\omega^{J(\frac{n}{2} - 1)})^{n-1} & (\omega^{J(\frac{n}{2} - 2)})^{n-1} & \cdots & (\omega^{J(0)})^{n-1} & (\omega^{J_*(\frac{n}{2} - 1)})^{n-1} & (\omega^{J_*(\frac{n}{2} - 2)})^{n-1} & \vdots  & (\omega^{J_*(0)})^{n-1}
\end{bmatrix}$}}

\textcolor{red}{ $\rhd$ 其中 $\omega = e^{i\pi/n} = \cos \left(\dfrac{\pi}{n}\right) + i\sin \left(\dfrac{\pi}{n}\right)$、$J(h) = 5^h \bmod 2n$，且 $J_*(h) = -5^h \bmod 2n$}

\small{\noindent $ = \begin{bmatrix}
1 & 1 & \cdots & 1 & 1 & 1 & \cdots & 1\\
(\omega^{J(\frac{n}{2} - 1)}) & (\omega^{J(\frac{n}{2} - 2)}) & \cdots & (\omega^{J(0)}) & (\overline\omega^{J(\frac{n}{2} - 1)}) & (\overline\omega^{J(\frac{n}{2} - 2)}) & \cdots & (\overline\omega^{J(0)})\\
(\omega^{J(\frac{n}{2} - 1)})^2 & (\omega^{J(\frac{n}{2} - 2)})^2 & \cdots & (\omega^{J(0)})^2 & (\overline\omega^{J(\frac{n}{2} - 1)})^2 & (\overline\omega^{J(\frac{n}{2} - 2)})^2 & \cdots & (\overline\omega^{J(0)})^2 \\
\vdots & \vdots & \ddots & \vdots & \vdots & \ddots & \vdots & \vdots \\
(\omega^{J(\frac{n}{2} - 1)})^{n-1} & (\omega^{J(\frac{n}{2} - 2)})^{n-1} & \cdots & (\omega^{J(0)})^{n-1} & (\overline\omega^{J(\frac{n}{2} - 1)})^{n-1} & (\overline\omega^{J(\frac{n}{2} - 2)})^{n-1} & \vdots  & (\overline\omega^{J(0)})^{n-1}
\end{bmatrix}$}

\textcolor{red}{ $\rhd$ 因为 $\omega^{-1} = e^{\frac{-i\pi}{n}} = \overline{e^{\frac{i\pi}{n}}} = \overline\omega$}

$ $

\item 将 $\vec{m}$ 转换为缩放后的整数向量 $\lceil \Delta\vec{m} \rfloor \approx \Delta \vec{m}$，其中 $\Delta$ 是大于 1 的尺度因子，使得 $\Delta m_i$ 在所有情形下，包括整个同态运算过程中，都不会超过 $q_0$ 对应范围的上下界（即满足 $0 \leq \Delta m_i < q_0$ 或 $-\dfrac{q_0}{2} \leq \Delta m_i < \dfrac{q_0}{2}$）。最终编码得到的明文多项式为 $\Delta M = \sum\limits_{i=0}^{n-1} \lceil \Delta m_i \rfloor X^i \text{ } \in \mathbb{Z}_q[X] / (X^n + 1)$。编码过程中的 $\lceil \Delta \vec{m} \rfloor$ 舍入会产生编码误差，因此 CKKS 是一种近似加密方案。

\end{enumerate}

\par\noindent\rule{\textwidth}{0.4pt}

\textbf{\underline{解码}：}由明文多项式 $\Delta M = \sum\limits_{i=0}^{n-1}\Delta m_iX^i$ 恢复 $\vec{m} = \dfrac{\Delta \vec{m}}{\Delta}$。然后，
计算 $\vec{v}_{'} = \hathat W^* \cdot \vec{m}$，其中：

\noindent $\hathat{W}^* = \begin{bmatrix}
1 & (\omega^{J(0)}) & (\omega^{J(0)})^2 & \cdots & (\omega^{J(0)})^{n-1}\\
1 & (\omega^{J(1)}) & (\omega^{J(1)})^2 & \cdots & (\omega^{J(1)})^{n-1}\\
1 & (\omega^{J(2)}) & (\omega^{J(2)})^2 & \cdots & (\omega^{J(2)})^{n-1}\\
\vdots & \vdots & \vdots & \ddots & \vdots \\
1 & (\omega^{J(\frac{n}{2}-1)}) & (\omega^{J(\frac{n}{2}-1)})^2 & \cdots & (\omega^{J(\frac{n}{2}-1)})^{n-1}\\
1 & (\omega^{J_*(0)}) & (\omega^{J_*(0)})^2 & \cdots & (\omega^{J_*(0)})^{n-1}\\
1 & (\omega^{J_*(1)}) & (\omega^{J_*(1)})^2 & \cdots & (\omega^{J_*(1)})^{n-1}\\
1 & (\omega^{J_*(2)}) & (\omega^{J_*(2)})^2 & \cdots & (\omega^{J_*(2)})^{n-1}\\
\vdots & \vdots & \vdots & \ddots & \vdots \\
1 & (\omega^{J_*(\frac{n}{2}-1)}) & (\omega^{J_*(\frac{n}{2}-1)})^2 & \cdots & (\omega^{J_*(\frac{n}{2}-1)})^{n-1}\\
\end{bmatrix}$

$ $

\noindent $ = \begin{bmatrix}
1 & (\omega^{J(0)}) & (\omega^{J(0)})^2 & \cdots & (\omega^{J(0)})^{n-1}\\
1 & (\omega^{J(1)}) & (\omega^{J(1)})^2 & \cdots & (\omega^{J(1)})^{n-1}\\
1 & (\omega^{J(2)}) & (\omega^{J(2)})^2 & \cdots & (\omega^{J(2)})^{n-1}\\
\vdots & \vdots & \vdots & \ddots & \vdots \\
1 & (\omega^{J(\frac{n}{2}-1)}) & (\omega^{J(\frac{n}{2}-1)})^2 & \cdots & (\omega^{J(\frac{n}{2}-1)})^{n-1}\\
1 & (\overline\omega^{J(0)}) & (\overline\omega^{J(0)})^2 & \cdots & (\overline\omega^{J(0)})^{n-1}\\
1 & (\overline\omega^{J(1)}) & (\overline\omega^{J(1)})^2 & \cdots & (\overline\omega^{J(1)})^{n-1}\\
1 & (\overline\omega^{J(2)}) & (\overline\omega^{J(2)})^2 & \cdots & (\overline\omega^{J(2)})^{n-1}\\
\vdots & \vdots & \vdots & \ddots & \vdots \\
1 & (\overline\omega^{J(\frac{n}{2}-1)}) & (\overline\omega^{J(\frac{n}{2}-1)})^2 & \cdots & (\overline\omega^{J(\frac{n}{2}-1)})^{n-1}\\
\end{bmatrix}$

$ $

$ $

最后，只提取\textit{正序}厄米向量 $\vec{v}_{'}$ 的前 $\dfrac{n}{2}$ 个元素，即可恢复输入向量 $\vec{v}$。

\end{tcolorbox}

\begin{comment}
\begin{table}[h] %usepackage{array}
\centering
\begin{tabular}{|>{\centering\arraybackslash}p{0.2\columnwidth}||>{\centering\arraybackslash}p{0.75\columnwidth}|}
\hline \hline
& \textbf{Encoded Polynomial $M(X) \in \mathcal{R}_{\langle n \rangle}$ over $X \in \mathbb{C}$ (complex numbers)} \\ \hline \hline
\textbf{Range of Input Slot Value}& Depends on the range of $M(X)$'s coefficient\\ \hline
\textbf{Range of $\bm{M(X)}$'s coefficients}&Any real number $m_i$ with the custom precision scaling factor $\Delta$ under the constraint: $0 \leq \Delta m_i \leq q-1$ (or $-\dfrac{q}{2} \leq \Delta m_i \leq \dfrac{q}{2} - 1$) \\\hline
%\textbf{Range of $\bm Y$}& $\left[0, \text{ } (1+i)\cdot {t}\cdot 2^{ \lceil \frac{n+1}{2} \rceil}\right]$\\
%& or $\left[-(1+i)\cdot \dfrac{t}{2}\cdot 2^{ \lceil \frac{n + 1}{2} \rceil}, \text{ } (1+i)\cdot \dfrac{t}{2}\cdot 2^{ \lceil \frac{n+1}{2} \rceil}\right]$ (a rough estimate)\\
\hline
\end{tabular}
\caption{The encoding capacity of $M(X) \in \mathcal{R}_{\langle n, t \rangle}$.}
\label{tab:polynomial-encoding-capacity-complex}
\end{table}

\para{CKKS's Encoding Range:} \autoref{tab:polynomial-encoding-capacity-complex} depicts the range of input values that can be encoded by the polynomials $M(X) \in \mathcal{R}_{\langle n \rangle}$ over $X \in \mathbb{C}$. As illustrated in \autoref{tab:polynomial-encoding-capacity-complex}, the encoding range of complex number inputs is determined by the size of the $n$ and $q$ parameters. Therefore, these two parameters should be chosen sufficiently large to encode all values that can be used by an application.
\end{comment}

\para{CKKS 的近似性质：} 令 $P=\lceil\Delta M\rfloor$ 表示对系数逐项舍入到最近整数，则 $P=\Delta M+E_{\mathrm{enc}}$，且 $\|E_{\mathrm{enc}}\|_\infty\leq 1/2$。由单位模根上的三角不等式，每个解码槽位的编码误差至多为 $n/(2\Delta)$。这是保守的系数范数界，应与 CKKS 分析采用的典范嵌入范数区分。加密和后续计算还会产生额外误差，近似解码不会自动将它们消除。


\subsubsection{示例}
\label{subsubsec:ckks-encoding-ex}


设输入复向量的维数为 $\dfrac{n}{2} = 2$，界定多项式的次数为 $n$ = 4，尺度因子为 $\Delta = 1024$。

所使用的 $n$ 维向量空间的基为

$\hathat{W}= \begin{bmatrix}
1 & 1 & 1 & 1\\
\omega^{J(1)} & \omega^{J(0)} & \overline{\omega}^{J(1)} & \overline{\omega}^{J(0)}\\
(\omega^{J(1)})^2 & ({\omega^{J(0)}})^2 & (\overline{\omega}^{J(1)})^2 & (\overline{\omega}^{J(0)})^2\\
(\omega^{J(1)})^3 & (\omega^{J(0)})^3 & (\overline{\omega}^{J(1)})^3 & (\overline{\omega}^{J(0)})^3\\
\end{bmatrix}$
$= \begin{bmatrix}
1 & 1 & 1 & 1\\
\omega^5 & \omega & \overline{\omega}^5 & \overline{\omega}\\
\omega^2 & \omega^2 & \overline{\omega^2} & \overline{\omega}^2\\
\omega^7 & \omega^3 & \overline{\omega}^7 & \overline{\omega}^3\\
\end{bmatrix}$

其中 $\omega = e^{i\pi/n} = \cos \left(\dfrac{\pi}{n}\right) + i\sin \left(\dfrac{\pi}{n}\right)$

$ $


在上述设置下，假设需要编码的输入复向量为 $\vec{v} = (1.1 + 4.3i, \text{ } 3.5 - 1.4i)$。

$ $

首先，构造正序厄米向量 $\vec{v}_{'} = (1.1 + 4.3i, \text{ } 3.5 - 1.4i, \text{ } 1.1 - 4.3i, \text{ } 3.5 + 1.4i)$。

$ $

接下来，应用如下变换，将复向量 $\vec{v}_{'}$ 转换为实数向量 $\vec{m}$：


$\vec{m} = \dfrac{\hathat{W} \cdot I_n^R \cdot \vec{v}_{'}}{n} = \dfrac{1}{4} \cdot \begin{bmatrix}
1 & 1 & 1 & 1\\
\omega^5 & \omega & \overline{\omega}^5 & \overline{\omega}\\
\omega^2 & \omega^2 & \overline{\omega}^2 & \overline{\omega}^2\\
\omega^7 & \omega^3 & \overline{\omega}^7 & \overline{\omega}^3\\
\end{bmatrix} \cdot
\begin{bmatrix}
0 & 0 & 0 & 1 \\
0 & 0 & 1 & 0 \\
0 & 1 & 0 & 0 \\
1 & 0 & 0 & 0
\end{bmatrix}
\cdot
\begin{bmatrix} 1.1 + 4.3i\\3.5 - 1.4i\\1.1 - 4.3i\\3.5 + 1.4i \end{bmatrix}$

$ $

$= \dfrac{\hathat W \cdot I_n^R \cdot \vec{v}_{'}}{n} = \dfrac{1}{4} \cdot\begin{bmatrix}
1 & 1 & 1 & 1\\
\overline{\omega} & \overline{\omega}^5 & \omega  & \omega^5\\
\overline{\omega}^2 & \overline{\omega}^2 & \omega^2 & \omega^2\\
\overline{\omega}^3 & \overline{\omega}^7 & \omega^3 & \omega^7\\
\end{bmatrix}
\cdot
\begin{bmatrix} 1.1 + 4.3i\\3.5 - 1.4i\\1.1 - 4.3i\\3.5 + 1.4i \end{bmatrix}$

$ $

$= \dfrac{1}{4} \cdot
\begin{bmatrix}
(1.1 + 4.3i) + (3.5 - 1.4i) + (1.1 - 4.3i) + (3.5 + 1.4i)\\
(1.1 + 4.3i)\overline{\omega} + (3.5 - 1.4i)\overline{\omega}^5 + (1.1 - 4.3i)\omega+ (3.5 + 1.4i)\omega^5 \\
(1.1 + 4.3i)\overline{\omega^2} + (3.5 - 1.4i)\overline{\omega^2} + (1.1 - 4.3i)\omega^2+ (3.5 + 1.4i)\omega^2 \\
(1.1 + 4.3i)\overline{\omega}^3 + (3.5 - 1.4i)\overline{\omega}^7 + (1.1 - 4.3i)\omega^3 + (3.5 + 1.4i)\omega^7 \\
\end{bmatrix}$

$ $

$= \dfrac{1}{4} \cdot
\begin{bmatrix}
(1.1 + 4.3i) + (3.5 - 1.4i) + (1.1 - 4.3i) + (3.5 + 1.4i)\\
(1.1 + 4.3i)\overline{\omega} + (3.5 - 1.4i)\overline{\omega}^5 + (1.1 - 4.3i)\omega+ (3.5 + 1.4i)\omega^5 \\
(1.1 + 4.3i)\overline{\omega^2} + (3.5 - 1.4i)\overline{\omega^2} + (1.1 - 4.3i)\omega^2+ (3.5 + 1.4i)\omega^2 \\
(1.1 + 4.3i)\overline{\omega}^3 + (3.5 - 1.4i)\overline{\omega}^7 + (1.1 - 4.3i)\omega^3 + (3.5 + 1.4i)\omega^7 \\
\end{bmatrix}$

$ $

$= \dfrac{1}{4} \cdot
\begin{bmatrix}
9.2\\
1.1(\overline{\omega} + \omega) + 4.3i(\overline{\omega} - \omega) +
3.5(\overline{\omega}^5 + \omega^5) - 1.4i(\overline{\omega}^5 - \omega^5)\\
1.1(\overline{\omega}^2 + {\omega}^2) + 4.3i(\overline{\omega}^2 - {\omega}^2) +
3.5(\overline{\omega}^2 + {\omega}^2) - 1.4i(\overline{\omega}^2 - {\omega}^2)\\
1.1(\overline{\omega}^3 + {\omega}^3) + 4.3i(\overline{\omega}^3 - {\omega}^3) +
3.5(\overline{\omega}^7 + {\omega}^7) - 1.4i(\overline{\omega}^7 - {\omega}^7)\\
\end{bmatrix}$

$ $

$= \dfrac{1}{4} \cdot
\begin{bmatrix}
9.2\\
1.1\left(2\cos\dfrac{\pi}{4}\right) - 4.3i\left(2i\sin\dfrac{\pi}{4}\right) + 3.5\left(2\cos\dfrac{5\pi}{4}\right) + 1.4i\left(2i\sin\dfrac{5\pi}{4}\right)\\
1.1\left(2\cos\dfrac{\pi}{2}\right) - 4.3i\left(2i\sin\dfrac{\pi}{2}\right) + 3.5\left(2\cos\dfrac{\pi}{2}\right) + 1.4i\left(2i\sin\dfrac{\pi}{2}\right)\\
1.1\left(2\cos\dfrac{3\pi}{4}\right) - 4.3i\left(2i\sin\dfrac{3\pi}{4}\right) + 3.5 \left(2\cos\dfrac{7\pi}{4}\right) + 1.4i\left(2i\sin\dfrac{7\pi}{4}\right)\\
\end{bmatrix}$

$ $

$= \dfrac{1}{4} \cdot
\begin{bmatrix}
9.2\\
1.1\left(2\dfrac{\sqrt{2}}{2}\right) + 4.3\left(2\dfrac{\sqrt{2}}{2}\right) + 3.5\left(-2\dfrac{\sqrt{2}}{2}\right) - 1.4\left(-2\dfrac{\sqrt{2}}{2}\right)\\
1.1(2\cdot 0) + 4.3(2\cdot1) + 3.5(2\cdot 0) - 1.4(2\cdot 1)\\
1.1\left(-2\dfrac{\sqrt{2}}{2}\right) + 4.3\left(2\dfrac{\sqrt{2}}{2}\right) + 3.5 \left(2\dfrac{\sqrt{2}}{2}\right) - 1.4\left(-2\dfrac{\sqrt{2}}{2}\right)\\
\end{bmatrix}$

$ $

$= 0.25 \cdot
\begin{bmatrix}
9.2\\
1.1\sqrt{2} + 4.3{\sqrt{2}} - 3.5\sqrt{2} + 1.4\sqrt{2}\\
1.1(0) + 4.3(2) - 3.5(0) - 1.4(2)\\
-1.1\sqrt{2} + 4.3\sqrt{2} + 3.5\sqrt{2} + 1.4\sqrt{2}\\
\end{bmatrix} =
\begin{bmatrix}
2.3\\
0.825\sqrt{2}\\
1.45\\
2.025\sqrt{2}\\
\end{bmatrix} \approx (2.3, \text{ } 1.1667, \text{ } 1.45, \text{ } 2.8638)$

$ $

将实数向量 $\vec{m}$ 乘以尺度因子 $\Delta$ 并舍入，得到缩放后的整数向量 $\Delta\vec{m}$：

$\Delta\vec{m} \approx \lceil \Delta \vec{m} \rfloor = \lceil 1024 \cdot (2.3, \text{ } 1.1667, \text{ } 1.45, \text{ } 2.8638) \rfloor = (2355, \text{ } 1195, \text{ } 1485, \text{ } 2933)$

$ $

最终，$\vec{v} = (1.1 + 4.3i, \text{ } 3.5 - 1.4i)$ 被编码为如下明文多项式 $M(X)$：

$\Delta M(X) = 2355 +  1195X + 1485X^2 + 2933X^3 \in \mathcal{R}_{\langle 4 \rangle}$

$ $


为了解码 $\vec{m}$，计算：

$\vec{v}_{'} = \dfrac{\hathat W^*\cdot\Delta\vec{m}}{\Delta} = \begin{bmatrix}
1,\omega,\omega^2, \omega^3\\
1,\omega^5, \omega^2,\omega^7\\
1,\overline{\omega}, \overline{\omega^2}, \overline{\omega^3}\\
1,\overline{\omega^5}, \overline{\omega^2}, \overline{\omega^7}
\end{bmatrix}$
$\cdot \begin{bmatrix}
2355\\1195\\1485\\2933
\end{bmatrix}\cdot \dfrac{1}{1024}$

$=
\begin{bmatrix}
1,\dfrac{\sqrt{2}}{2} + \dfrac{i\sqrt{2}}{2},i, -\dfrac{\sqrt{2}}{2} + \dfrac{i\sqrt{2}}{2}\\
1,-\dfrac{\sqrt{2}}{2} - \dfrac{i\sqrt{2}}{2}, i,\dfrac{\sqrt{2}}{2} - \dfrac{i\sqrt{2}}{2}\\
1,\dfrac{\sqrt{2}}{2} - \dfrac{i\sqrt{2}}{2}, -i, -\dfrac{\sqrt{2}}{2} - \dfrac{i\sqrt{2}}{2}\\
1,-\dfrac{\sqrt{2}}{2} + \dfrac{i\sqrt{2}}{2}, -i, \dfrac{\sqrt{2}}{2} + \dfrac{i\sqrt{2}}{2}
\end{bmatrix}$
$\cdot \begin{bmatrix}
2.2998046875\\1.1669921875\\1.4501953125\\2.8642578125
\end{bmatrix}$

$ $

$ \approx (1.0997+4.3007i, \text{ } 3.5000-1.4003i, \text{ } 1.0997-4.3007i, \text{ } 3.5000+1.4003i)$

$ $

提取厄米向量 $\vec{v}_{'}$ 的前 $\dfrac{n}{2} = 2$ 个元素，即可恢复输入向量：

$(1.0997+4.3007i, \text{ } 3.5000-1.4003i)$

$\approx (1.1 + 4.3i, \text{ } 3.5 - 1.4i) = \vec{v}$ \textcolor{red}{  $\rhd$ 原始输入向量}

$ $

对放大后的逆嵌入作整数舍入，产生了原始槽位与解码槽位之间的小差异。所展示的中间小数系数是近似值，整数编码应根据含 $\sqrt2$ 的精确表达式计算。

$ $

\para{源代码：}运行以下文件可执行 CKKS 编码示例：\href{https://github.com/wang2000star/fhe-textbook-cn/blob/chinese-edition/source%20code/ckks.py}{\underline{此 Python 脚本}}.



\subsection{加密与解密}
\label{subsec:ckks-enc-dec}

CKKS 加密的是 $P=\lceil\Delta M\rfloor\in R$，而不是直接加密任意实数系数。下面给出对称密钥 RLWE 形式；公钥 CKKS 加密使用公开的 RLWE 样本，并具有相应的组合误差。

\begin{tcolorbox}[title={\textbf{\tboxlabel{\ref*{subsec:ckks-enc-dec}} CKKS 加密与解密}}]
选择小秘密 $S\in R$，均匀采样 $A\leftarrow R_q$ 及小整数误差 $E_0$，输出
\[
B=[-AS+P+E_0]_q,\qquad \textsf{ct}=(A,B)\in R_q^2.
\]
定义原始解密相位为
\[
\mathsf{Phase}_S(\textsf{ct})=[B+AS]_q=P+E_0=\Delta M+E,
\qquad E=E_{\mathrm{enc}}+E_0,
\]
其成立条件为 $\|\Delta M+E\|_\infty<q/2$。这里 $[\cdot]_q$ 选择中心化系数代表元。解码时，在选定的根上求相位的值，再除以 $\Delta$：
\[
\widehat v_j=\Delta^{-1}\mathsf{Phase}_S(\textsf{ct})(\omega^{J(j)})
=v_j+\Delta^{-1}E(\omega^{J(j)}).
\]
因此 $|\widehat v_j-v_j|\leq\|E\|_{\mathrm{can},\infty}/\Delta\leq n\|E\|_\infty/\Delta$。仅有 $|e_i|<\Delta/2$ 并不能推出槽位精度为 $1/\Delta$。舍入到 $1/\Delta$ 的整数倍也不能消除一般的 CKKS 噪声。
\end{tcolorbox}

下文在代数推导中使用 $\textsf{RLWE}^{-1}$ 时，均指这一原始相位；除以所记录的尺度和槽位解码是独立操作。$\textsf{RLWE}_{S,\sigma}(P)$ 表示加密关系，并不声称计算后的密文仍具有原来的独立误差分布 $\chi_\sigma$。

\subsection{密文与密文相加}
\label{subsec:ckks-add-cipher}

CKKS 与 BFV 一样逐分量相加密文（总结~\ref*{subsec:bfv-add-cipher}、\autoref{subsec:bfv-add-cipher}）。密文的密钥、模数和尺度必须相容。结果的相位误差为 $E^{\langle1\rangle}+E^{\langle2\rangle}$，相位之和必须满足不回绕界。




\begin{tcolorbox}[title={\textbf{\tboxlabel{\ref*{subsec:ckks-add-cipher}} CKKS 密文与密文相加}}]
$\textsf{RLWE}_{S, \sigma}(\Delta M^{\langle 1 \rangle} ) + \textsf{RLWE}_{S, \sigma}(\Delta M^{\langle 2 \rangle} ) $

$ = ( A^{\langle 1 \rangle}, \text{ } B^{\langle 1 \rangle}) + (A^{\langle 2 \rangle}, \text{ } B^{\langle 2 \rangle}) $

$ = ( A^{\langle 1 \rangle} + A^{\langle 2 \rangle}, \text{ } B^{\langle 1 \rangle} + B^{\langle 2 \rangle} ) $

$= \textsf{RLWE}_{S, \sigma}(\Delta(M^{\langle 1 \rangle} + M^{\langle 2 \rangle}) )$
\end{tcolorbox}


\subsection{密文与明文相加}
\label{subsec:ckks-add-plain}

CKKS 与 BFV 一样将明文加到 $B$ 分量（总结~\ref*{subsec:bfv-add-plain}、\autoref{subsec:bfv-add-plain}）。应按密文的实际尺度 $\Delta$ 编码明文，并将其编码误差计入误差预算。下述简写中的 $\Delta\Lambda$ 表示舍入后的整数编码。



\begin{tcolorbox}[title={\textbf{\tboxlabel{\ref*{subsec:ckks-add-plain}} CKKS 密文与明文相加}}]
$\textsf{RLWE}_{S, \sigma}(\Delta M) + \Delta\Lambda $

$=  (A, \text{ } B) + \Delta\Lambda$

$=  (A, \text{ } B + \Delta\cdot\Lambda)$

$= \textsf{RLWE}_{S, \sigma}(\Delta (M + \Lambda) )$
\end{tcolorbox}




\subsection{密文与密文相乘}
\label{subsec:ckks-mult-cipher}

CKKS 乘法先在当前模数下对密文作张量积，然后通常进行重线性化和重缩放。张量积将输入尺度 $\Delta_1,\Delta_2$ 变为 $\Delta_1\Delta_2$；除以 $w_l=q_l/q_{l-1}$ 的重缩放将其变为 $\Delta_1\Delta_2/w_l$，并消耗一层。下述推导采用定理~\ref*{subsec:glwe-alternative}（\autoref{subsec:glwe-alternative}）中的符号约定 $B=-AS+\Delta M+E\pmod q$。我们先推导包含误差的精确相位，再说明两种重线性化方法及重缩放。


\subsubsection{合成密文的推导}
\label{subsubsec:ckks-mult-cipher-relation}

对于模数和秘密密钥相同的两个密文，设其相位分别为 $\Delta M^{\langle 1\rangle}+E^{\langle 1\rangle}$ 和 $\Delta M^{\langle 2\rangle}+E^{\langle 2\rangle}$。在模 $q$ 下定义
\[
D_0=B^{\langle 1\rangle}B^{\langle 2\rangle},\quad
D_1=A^{\langle 1\rangle}B^{\langle 2\rangle}+A^{\langle 2\rangle}B^{\langle 1\rangle},\quad
D_2=A^{\langle 1\rangle}A^{\langle 2\rangle}.
\]
精确的相位恒等式为
\[
D_0+D_1S+D_2S^2\equiv\Delta^2M^{\langle 1\rangle}M^{\langle 2\rangle}+F\pmod q,
\]
\[
F=\Delta M^{\langle 1\rangle}E^{\langle 2\rangle}
+\Delta M^{\langle 2\rangle}E^{\langle 1\rangle}
+E^{\langle 1\rangle}E^{\langle 2\rangle}.
\]
CKKS 解码会保留这些误差项。重线性化构造 $C_\alpha=(D_1,D_0)$，以及相位为 $D_2S^2+E_{\mathrm{ks}}$ 的多项式对 $C_\beta$。于是 $C_\alpha+C_\beta\in R_q^2$ 的相位为 $\Delta^2M^{\langle 1\rangle}M^{\langle 2\rangle}+F+E_{\mathrm{ks}}$。这并未将尺度恢复为 $\Delta$；恢复尺度是随后重缩放步骤的作用。如果不进行重线性化，密文长度会随关于 $S$ 的次数增长；反复平方会在每一步将该次数翻倍。

\subsubsection{重线性化方法 1：密文分解}
\label{subsubsec:relinearization-gadget-decomposition}

对于任意 $q$，使用显式的整数 gadget：选择整数 $\beta\geq 2$，令 $l=\lceil\log_\beta q\rceil$，并取 $g_i=\beta^{i-1}$，其中 $1\leq i\leq l$。将 $D_2$ 各系数的无符号代表元分解为数位 $d_i\in\{0,\ldots,\beta-1\}$，使
\[
D_2\equiv\sum_{i=1}^l d_i g_i\pmod q.
\]
这里的升序 gadget 与降序权重 $q/\beta^i$ 是不同约定，同一个数位向量不能同时配用这两套权重。生成求值密钥对 $K_i=(A_i,B_i)$，使其满足
\[
B_i+A_iS\equiv g_iS^2+E_i\pmod q.
\]
于是
\[
C_\beta=\sum_{i=1}^l d_iK_i,\qquad
\mathsf{Phase}_S(C_\beta)\equiv D_2S^2+\sum_{i=1}^l d_iE_i\pmod q.
\]
对于负循环系数范数，有 $\|\sum_i d_iE_i\|_\infty\leq n(\beta-1)\sum_i\|E_i\|_\infty$。使用平衡数位或截断分解时，必须采用相应的数位界与残差界。求值密钥加密的是 $S^2$ 的 gadget 倍数，不附加 CKKS 消息尺度 $\Delta$。

\subsubsection{重线性化方法 2：密文模数切换}
\label{subsubsec:relinearization-modulus-switch}

另一种重线性化方法使用辅助正整数 $g$ 和模 $gq$ 下的求值密钥。应根据密钥切换的误差预算选择 $g$，它不必随机，也不必是二次幂。在 $R_{gq}$ 中均匀采样 $A'$，并采样小整数误差 $E'$，令
\[
\mathit{evk}_g=(A',B'),\qquad B'=[-A'S+gS^2+E']_{gq}.
\]
将 $D_2\in R_q$ 的系数提升为固定的中心化整数。计算多项式对 $Z=[D_2\mathit{evk}_g]_{gq}$，并令
\[
C_\beta=\left[\left\lceil Z/g\right\rfloor\right]_q.
\]
记 $R_A=\lceil Z_A/g\rfloor-Z_A/g$、$R_B=\lceil Z_B/g\rfloor-Z_B/g$，其每个系数的绝对值均不超过 $1/2$。在模 $q$ 下，输出相位满足
\[
\mathsf{Phase}_S(C_\beta)\equiv D_2S^2+D_2E'/g+R_B+R_AS\pmod q.
\]
因此，新增误差的系数范数满足
\[
\|E_{\mathrm{ks}}\|_\infty\leq
\frac{n\|D_2\|_\infty\|E'\|_\infty}{g}
+\frac{1+\|S\|_1}{2}.
\]
对 $Z$ 取模产生的 $gq$ 倍数，在除法后变为 $q$ 的倍数，因此在模 $q$ 下消失。舍入不对加法分配，也不能与乘以 $S$ 交换；其贡献应显式写成 $R_B+R_AS$。较大的 $g$ 会抑制求值密钥误差，但不能消除此舍入项。$g$ 的选择还会影响扩展模数求值密钥的安全参数。


\subsubsection{重缩放}
\label{subsubsec:ckks-mult-cipher-rescale}

设 $q_l=w_lq_{l-1}$，重缩放前的密文相位为 $\Delta_{\mathrm{in}}M+F$。重缩放将中心化密文代表元逐系数除以 $w_l$，再舍入到最近整数：
\[
\widehat A=[\lceil A/w_l\rfloor]_{q_{l-1}},\qquad
\widehat B=[\lceil B/w_l\rfloor]_{q_{l-1}}.
\]
令 $R_A=\lceil A/w_l\rfloor-A/w_l$、$R_B=\lceil B/w_l\rfloor-B/w_l$，则
\[
\widehat B+\widehat AS\equiv
\frac{\Delta_{\mathrm{in}}}{w_l}M+
\frac{F}{w_l}+R_B+R_AS\pmod {q_{l-1}}.
\]
因此实际输出尺度为 $\Delta_{\mathrm{out}}=\Delta_{\mathrm{in}}/w_l$，尺度元数据本身不作舍入。对于尺度分别为 $\Delta_1,\Delta_2$ 的密文乘积，有 $\Delta_{\mathrm{out}}=\Delta_1\Delta_2/w_l$。若 $\Delta_1=\Delta_2=\Delta$ 且 $w_l\approx\Delta$，结果接近 $\Delta$，但一般不与之相等。将尺度重新标为精确的 $\Delta$ 会引入依赖消息的误差，必须另行分析精度。

新的相位误差为 $\widehat F=F/w_l+R_B+R_AS$，不是逐系数的 $\lceil F/w_l\rfloor$。确定性界为
\[
\|\widehat F\|_\infty\leq\|F\|_\infty/w_l+(1+\|S\|_1)/2.
\]
正确解释结果还要求完整的输出相位位于模 $q_{l-1}$ 的中心化区间内。

对于 RNS 链 $q_l=\prod_{i=0}^l w_i$，应根据算术效率和目标尺度安排选择素数。RNS 改变的是表示形式，而不是 $q_l$ 的总信息量。RLWE 安全性取决于维数、总模数、秘密分布、误差分布和求值密钥参数，而不是仅由最小及最大的 RNS 素数决定。

\para{加密时的初始尺度因子 $\Delta$：} 应根据目标精度和可用模数空间选择 $\Delta$，并使 $q_0$ 足以容纳预期最终消息及误差。令 $\Delta$ 接近重缩放素数通常较为方便，但并非必须取 $\Delta=w_L$。

\para{噪声增长：} 乘法会传播消息与误差的交叉项及误差乘积，而重缩放会新增舍入误差。界取决于消息与误差的大小；并非所有输入的误差都会增长，例如乘以零时就未必如此。单独重缩放不能恢复已经损失的数值精度。

\subsubsection{BFV 与 CKKS 自举的比较}
\label{subsubsec:bfv-bootstrapping-ckks-comparison}

两种方案都通过求值适当的解密相关函数来刷新密文，但具体步骤与明文含义不同。在下文的 CKKS 流程中，提升密文系数会在较大模数下暴露原始相位 $\Delta M+E+q_0K$，无需另行计算包含加密秘密的内积。线性的系数与槽位变换以及近似模约减函数会近似消除 $q_0K$ 的贡献。BFV 自举则求值有限明文空间上的精确解密函数，并须满足相应参数和数位提取条件。

\para{乘法层级：} 即使保持固定密文模数，BFV 仍有有限的噪声预算。在 CKKS 标准的乘法、重线性化、重缩放流程中，重缩放会消耗一层模数链。到达 $q_0$ 意味着不能再按该安排进行重缩放，而不是在代数上禁止一切乘法。自举可补充可用层级，同时也会消耗若干层（\autoref{subsec:ckks-bootstrapping}）。

\para{噪声处理的局限：} 重缩放将原有相位误差和模数除以同一因子，并增加舍入误差。因此，它本身不会改善原有误差与模数的比值，也不能恢复损失的数值精度。CKKS 自举扩大可用模数，但会保留已有数值误差，并增加自身的近似与求值误差。


\subsubsection{总结}
\label{subsubsec:ckks-mult-cipher-summary}

\begin{tcolorbox}[title={\textbf{\tboxlabel{\ref*{subsubsec:ckks-mult-cipher-summary}} CKKS 密文与密文乘法}}]
对于密钥和模数 $q_l$ 相同、尺度分别为 $\Delta_1,\Delta_2$ 的密文：
\begin{enumerate}
\item 在模 $q_l$ 下计算 $D_0=B_1B_2$、$D_1=A_1B_2+A_2B_1$ 和 $D_2=A_1A_2$。它们的相位是两个输入相位的乘积。
\item 使用前述任一种方法重线性化，得到 $(A,B)=(D_1,D_0)+C_\beta$，其相位为 $\Delta_1\Delta_2 M_1M_2+F$；$F$ 包含乘法与密钥切换误差。
\item 若计算安排要求重缩放，令 $q_{l-1}=q_l/w_l$，并计算
\[
(\widehat A,\widehat B)=[(\lceil A/w_l\rfloor,\lceil B/w_l\rfloor)]_{q_{l-1}},
\quad \Delta_{\mathrm{out}}=\Delta_1\Delta_2/w_l.
\]
输出误差为 $F/w_l+R_B+R_AS$，并须满足上述不回绕与精度界。
\end{enumerate}
\para{交换\textsf{重线性化}与\textsf{重缩放}的顺序：} 可以先重缩放三分量密文，其舍入误差为 $R_0+R_1S+R_2S^2$，再使用适用于较低模数的密钥重线性化。两种安排并不是相同的密文计算，其误差和成本界也不同。
\end{tcolorbox}


\subsection{密文与明文相乘}
\label{subsec:ckks-mult-plain}

以选定的尺度 $\Delta_\Lambda$ 编码明文乘数 $\Lambda$，得到 $P_\Lambda=\lceil\Delta_\Lambda\Lambda\rfloor$。对于尺度为 $\Delta$ 的密文 $(A,B)$，计算
\[
(A',B')=(AP_\Lambda,BP_\Lambda)\pmod q,
\qquad \Delta'=\Delta\Delta_\Lambda.
\]
相位误差包含输入密文误差乘以 $P_\Lambda$ 的贡献，以及明文编码误差的贡献。若需要除以 $w_l$ 重缩放，则使用 \autoref{subsubsec:ckks-mult-cipher-summary} 的流程，输出尺度为 $\Delta\Delta_\Lambda/w_l$，并增加相应的舍入误差。这种常见安排消耗一层，但并非每次明文乘法都必须如此：精确表示的小整数常量可以采用 $\Delta_\Lambda=1$；模数和精度预算允许时，也可以延后重缩放。


\subsection{\textsf{ModDrop}}
\label{subsec:ckks-moddrop}

对于 $q_{l-1}\mid q_l$，定义
\[
(A',B')=([A]_{q_{l-1}},[B]_{q_{l-1}}).
\]
若 $A'=A-K_Aq_{l-1}$、$B'=B-K_Bq_{l-1}$，则
\[
A'S+B'=AS+B-(K_AS+K_B)q_{l-1}.
\]
因此新相位与原相位模 $q_{l-1}$ 同余。若 $\|\Delta M+E\|_\infty<q_{l-1}/2$，其中心化代表元仍精确等于 $\Delta M+E$；否则相位可能回绕，从而改变解码值。

\begin{tcolorbox}[title={\textbf{\tboxlabel{\ref*{subsec:ckks-moddrop}} CKKS \textsf{ModDrop}}}]
将每个密文分量对当前模数的一个因子取模。只有完整相位落在新的中心化区间内，所记录的尺度和已有相位误差才保持不变；可用模数空间会减少。
\end{tcolorbox}

\subsubsection{模数切换与 \textsf{ModDrop} 的区别}
\label{subsubsec:ckks-moddrop-vs-modswitch}

两种操作都会降低密文模数。重缩放还将记录的尺度和原有误差除以 $q_l/q_{l-1}$，并引入舍入误差。\textsf{ModDrop} 保持尺度不变，但只有满足其不回绕条件时才保持相位。它也可用于加法或乘法前对齐不同层级的密文，并非只能在不再进行乘法时使用。

\subsection{同态密钥切换}
\label{subsec:ckks-key-switching}

选取 \autoref{subsubsec:relinearization-gadget-decomposition} 中的精确整数 gadget $g_i$，并分解 $A\equiv\sum_i d_ig_i\pmod q$。新秘密密钥 $S'$ 下的求值密钥 $(a_i,b_i)$ 满足 $b_i+a_iS'\equiv g_iS+E_i\pmod q$。于是
\begin{tcolorbox}[title={\textbf{\tboxlabel{\ref*{subsec:ckks-key-switching}} CKKS 密钥切换}}]
\[
(A',B')=(0,B)+\sum_i d_i(a_i,b_i),\qquad
B'+A'S'\equiv B+AS+\sum_i d_iE_i\pmod q.
\]
\end{tcolorbox}
尺度不变，误差增加所列密钥切换项。正确解释要求新相位适合该模数。这与 \autoref{subsec:bfv-key-switching} 使用相同的基础 gadget 技术，但采用 CKKS 的近似解码含义。


\subsection{输入向量槽位的同态旋转}
\label{subsec:ckks-rotation}

BFV 与 CKKS 槽位旋转采用相同的根指数置换，但 BFV 在其明文环中的根上求值，CKKS 则在复数根上求值。在 \autoref{subsec:ckks-encoding-decoding} 的编码约定中，逆变换为 $n^{-1}\hathat W I_n^R$，解码变换为 $\hathat W^*$，矩阵顺序不可颠倒。

对于 $J(j)=5^j\bmod2n$，自同构 $M(X)\mapsto M(X^{J(h)})$ 将槽位 $j$ 变为原来的槽位 $j+h\bmod(n/2)$，因为 $5^j5^h=5^{j+h}$。在 $J_*(j)=-5^j\bmod2n$ 的约定下，共轭半部向相同方向旋转。
\begin{tcolorbox}[title={\textbf{\tboxlabel{\ref*{subsec:ckks-rotation}} CKKS 输入向量槽位的同态旋转}}]
\begin{enumerate}
\item 对 $A$ 与 $B$ 都施加 $X\mapsto X^{J(h)}$，并模 $X^n+1$ 约减。秘密密钥变为 $S(X^{J(h)})$。
\item 使用 \autoref{subsec:ckks-key-switching} 中的精确 gadget 密钥切换，从 $S(X^{J(h)})$ 切换到 $S(X)$。将新增密钥切换误差计入输出相位。
\end{enumerate}
结果将 $n/2$ 个独立复数槽位向左旋转 $h$ 个位置。
\end{tcolorbox}

\para{半部槽位内的旋转：} 单个 $5^h$ 自同构分别旋转两个轨道，但并不禁止使用掩码与若干自同构组合构造其他置换。CKKS 的后半部由共轭关系确定，而完全分裂的 BFV 打包允许独立选择全部 $n$ 个槽位。


\subsubsection{示例}
\label{subsubsec:ckks-rotation-ex}

本小节给出以下两个示例：
\begin{enumerate}
\item 基于更新后的编码与解码矩阵 $\hathat W$ 和 $\hathat W^*$，将输入向量 $\vec{v}$ 编码为明文多项式 $M(X)$
\item 将编码后的明文 $M(X)$ 更新为 $M(X^{J(h)})$，使输入向量 $\vec{v}$ 的所有元素向左旋转 $h$ 个位置
\end{enumerate}

$ $

仍使用 \autoref{subsubsec:ckks-encoding-ex} 中相同的输入向量 $\vec{v}$：$\vec{v} = (1.1 + 4.3i, 3.5 - 1.4i)$。

回顾一下，$\vec{v}$ 编码得到的明文多项式如下：

$\Delta M(X) = 2355 + 1195X + 1485X^2 + 2933X^3 \in \mathcal{R}_{\langle 4 \rangle}=\mathbb{R}[X]/(X^4+1)$


假设希望将输入向量 $\vec{v}$ 向左旋转 1 个位置：

$\vec{v}^{\langle h=1 \rangle} = (3.5 - 1.4i, 1.1 + 4.3i)$

$ $

因此，将 $\Delta M(X)$ 更新为 $\Delta M(X^{J(1)})$：

$\Delta M(X^{J(1)}) = \Delta M(X^{5}) = 2355 + 1195(X^5) + 1485(X^5)^2 + 2933(X^5)^3$

$= 2355 + 1195X^5 + 1485X^{10} + 2933X^{15}$

$= 2355 + 1195X\cdot(-1) + 1485X^{2}\cdot(-1)\cdot(-1) + 2933X^{3}\cdot(-1)\cdot(-1)\cdot(-1)$

$= 2355 - 1195X + 1485X^2 - 2933X^3$

$ $

旋转后的\textit{正序}厄米输入向量计算如下：

$\dfrac{\hathat W^* \cdot \Delta \vec{m}}{\Delta} = \begin{bmatrix}
1,\omega,\omega^2, \omega^3\\
1,\omega^5, \omega^{10},\omega^{15}\\
1,\overline{\omega}, \overline{\omega^2}, \overline{\omega^3}\\
1,\overline{\omega^5}, \overline{\omega^{10}}, \overline{\omega^{15}}
\end{bmatrix} \cdot \begin{bmatrix}
2355\\- 1195\\1485\\-2933
\end{bmatrix}\cdot \dfrac{1}{1024}$

$= \dfrac{\hathat W^* \cdot \Delta \vec{m}}{\Delta} = \begin{bmatrix}
1,\omega,\omega^2, \omega^3\\
1,\omega^5, \omega^{2},\omega^{7}\\
1,\overline{\omega}, \overline{\omega^2}, \overline{\omega^3}\\
1,\overline{\omega^5}, \overline{\omega^{2}}, \overline{\omega^{7}}
\end{bmatrix} \cdot \begin{bmatrix}
2355\\- 1195\\1485\\-2933
\end{bmatrix}\cdot \dfrac{1}{1024}$

$= \dfrac{\hathat W^* \cdot \Delta \vec{m}}{\Delta} = \begin{bmatrix}
1,\omega,\omega^2, \omega^3\\
1,\omega^5, \omega^{2},\omega^{7}\\
1,\omega^7, \omega^6, \omega^5\\
1,\omega^3, \omega^{6}, \omega
\end{bmatrix} \cdot \begin{bmatrix}
2355\\- 1195\\1485\\-2933
\end{bmatrix}\cdot \dfrac{1}{1024}$

$=
\begin{bmatrix}
1,\dfrac{\sqrt{2}}{2} + \dfrac{i\sqrt{2}}{2},i, -\dfrac{\sqrt{2}}{2} + \dfrac{i\sqrt{2}}{2}\\
1,-\dfrac{\sqrt{2}}{2} - \dfrac{i\sqrt{2}}{2}, i,\dfrac{\sqrt{2}}{2} - \dfrac{i\sqrt{2}}{2}\\
1,\dfrac{\sqrt{2}}{2} - \dfrac{i\sqrt{2}}{2}, -i, -\dfrac{\sqrt{2}}{2} - \dfrac{i\sqrt{2}}{2}\\
1,-\dfrac{\sqrt{2}}{2} + \dfrac{i\sqrt{2}}{2}, -i, \dfrac{\sqrt{2}}{2} + \dfrac{i\sqrt{2}}{2}
\end{bmatrix}$
$\cdot \begin{bmatrix}
2.2998046875\\-1.1669921875\\1.4501953125\\-2.8642578125
\end{bmatrix}$

$ $

$ \approx (3.500 - 1.4003i, \text{ } 1.0997 + 4.3007i,\text{ } 3.500 + 1.4003i, \text{ } 1.0997 - 4.3007i)$

$ $

提取上述厄米向量的前 $\dfrac{n}{2} = 2$ 个元素，即可恢复输入向量：

$(3.500 - 1.4003i, \text{ } 1.0997 + 4.3007i)$

$ \approx (3.5 - 1.4i, \text{ } 1.1 + 4.3i)$
$ = \vec{v}^{\langle h=1 \rangle}$ \textcolor{red}{  $\rhd$ 原始输入向量 $\vec{v}$ 向左旋转 1 个位置}



$ $

实际求值时，对密文分量施加自同构，得到
\[
(A(X^{J(1)}),B(X^{J(1)})).
\]
此密文的秘密密钥为 $S(X^{J(1)})$。在使用原密钥解密前，应先将其密钥切换回 $S(X)$。解码结果近似于旋转后的向量 $\vec v^{\langle h=1\rangle}$，误差包含原编码误差及新增的密钥切换误差。

$ $


\para{源代码：}运行以下文件可执行 CKKS 输入向量槽位同态旋转的示例：\href{https://github.com/wang2000star/fhe-textbook-cn/blob/chinese-edition/source%20code/ckks.py}{\underline{此 Python 脚本}}.





\subsection{对 CKKS 编码的进一步思考}
\label{subsec:ckks-encoding-contemplate}

\para{CKKS 编码为何使用 $\bm{(\mu=2n)}$ 次分圆多项式：} 在 $X^n+1$ 的 $n$ 个不同根上求值，会得到可逆的范德蒙德矩阵。在域上，任意一组互异求值点都给出可逆范德蒙德矩阵；单位根还使求值与模 $X^n+1$ 的乘法相容，提供共轭配对，并支持高效傅里叶变换。二次幂分圆环是一种有用的实现选择，而不是因为只有它们的范德蒙德矩阵可逆。

采用槽位顺序 $\omega^{5^j}$ 时，施加 $X\mapsto X^{5^h}$ 会置换求值矩阵（即解码矩阵）的行。逆变换才是编码矩阵，不能将两者的方向混淆。共轭配对将独立复数坐标数降为 $n/2$，同时保留 $n$ 个实系数坐标。


\subsection{同态共轭}
\label{subsec:ckks-conjugation}

对于实系数多项式 $M$ 和 $\zeta_j=\omega^{J(j)}$，共轭满足
\[
\overline{M(\zeta_j)}=M(\overline{\zeta_j})=M(\zeta_j^{-1}).
\]
因此，正序完整向量由 $(\vec v,\overline{\vec v})$ 变为 $(\overline{\vec v},\vec v)$。应使用与编码一致的选定根 $J(j)=5^j\bmod 2n$，不能将其顺序替换为 $1,3,\ldots,n-1$。

\begin{tcolorbox}[title={\textbf{\tboxlabel{\ref*{subsec:ckks-conjugation}} CKKS 同态共轭}}]
\begin{enumerate}
\item 对密文的两个分量施加环自同构 $X\mapsto X^{-1}=X^{2n-1}$，并对 $X^n+1$ 取模。
\item 从 $S(X^{-1})$ 密钥切换到 $S(X)$，同时计入新增密钥切换误差。
\end{enumerate}
\end{tcolorbox}

\subsection{密文的稀疏打包}
\label{subsec:ckks-sparse-packing}

在 \autoref{subsec:ckks-encoding-decoding} 中，我们学习了 CKKS 编码方案：将具有 $\dfrac{n}{2}$ 个槽位的输入向量（即 $\dfrac{n}{2}$ 槽位输入向量）编码为 $(n-1)$ 次多项式。虽然多项式环的次数 $n$ 作为安全参数在 CKKS 设置阶段便已固定，但某些应用需要的输入向量槽位数可能少于 $\dfrac{n}{2}$。假设在 $\frac{n}{2}$ 个槽位中仅需使用 $\frac{n'}{2}$ 个槽位，其中 $n'\geq2$ 且整除 $n$，则相应的输入向量与编码后的多项式具有以下特殊性质：


\begin{tcolorbox}[title={\textbf{\tboxlabel{\ref*{subsec:ckks-sparse-packing}} CKKS 的稀疏打包多项式与密文}}]

假设一个 $\frac{n}{2}$ 槽位输入向量被编码为 $\mathbb{R}[X] / (X^n + 1)$ 中的多项式，并设 $n'\geq2$ 且整除 $n$。定义多项式 $M(X) \in \mathbb{R}[X] / (X^n + 1)$：仅指数为 $\frac{n}{n'}$ 的倍数的项可以具有非零系数，其余项的系数均为零（即 $M(X) = c_0 + c_{\frac{n}{n'}}X^{\frac{n}{n'}} + c_{\frac{2n}{n'}}X^{\frac{2n}{n'}} + \cdots + c_{n - \frac{n}{n'}}X^{n - \frac{n}{n'}}$）。于是可将 $M(X)$ 写为某个 $M_Y(Y) \in \mathbb{R}[Y] / (Y^{n'} + 1)$，其中 $Y = X^{\frac{n}{n'}}$，因此 $M(X) = M_Y(X^{\frac{n}{n'}})$。

$ $

那么，以下结论成立：

\begin{enumerate}
    \item 每个 $M_Y(Y) \in \mathbb{R}[Y]/(Y^{n'} + 1)$ 都同构地映射（即解码）为某个 $\dfrac{n}{2}$ 槽位输入向量，该向量由连续 $\dfrac{n'}{2}$ 个槽位值重复 $\dfrac{n}{n'}$ 次构成。
    \item 反过来，如果一个 $\dfrac{n}{2}$ 槽位输入向量由前 $\dfrac{n'}{2}$ 个连续槽位值重复 $\dfrac{n}{n'}$ 次构成，那么该向量将编码为某个 $M_Y(Y) \in \mathbb{R}[Y] / (Y^{n'} + 1)$，即 $\mathbb{R}[X] / (X^n + 1)$ 中指数不是 $\dfrac{n}{n'}$ 的倍数的各项系数均为零的多项式。
\end{enumerate}

\end{tcolorbox}

假设 $n'\geq2$ 且整除二次幂 $n$。下面证明正向包含关系。反向结论还需要注意：子环 $\mathbb{R}[Y]/(Y^{n'}+1)$ 与重复槽位向量空间的实维数都为 $n'$。单射的典范嵌入将前者映入后者，因此维数相等保证它满射到该子空间。仅指出整个编码映射是同构，而不论证这些子空间，尚不足以证明反向结论。

\subsubsection{正向证明：稀疏打包密文的解码}

下面证明，对于每个 $M_Y(Y) \in \mathbb{R}[Y] / (Y^{n'} + 1)$（即 $\mathbb{R}[X] / (X^n + 1)$ 中仅指数为 $\dfrac{n}{n'}$ 的倍数的项可能具有非零系数、其余项系数均为零的多项式），其解码结果是某个 $\dfrac{n}{2}$ 槽位输入向量，由前 $\dfrac{n'}{2}$ 个连续槽位值重复 $\dfrac{n}{n'}$ 次构成。

要将 $M(X)$ 解码为输入向量，应从 $X^n + 1$ 的 $n$ 个不同本原 $(\mu=2n)$ 次单位根中选取 $\frac{n}{2}$ 个，对 $M(X)$ 求值。求值向量为：

$\bm{(} \text{ }
M(\omega^{J(0)}), \text{ } M(\omega^{J(1)}), \cdots,  M(\omega^{J(\frac{n}{2}-1)})\bm{)}$

\textcolor{red}{ $\rhd$ 其中 $\omega=e^{\frac{i\pi}{n}}$ 是本原（$\mu=2n$）次单位根的基值与生成元}

$ $

由于 $M(X) = M_Y(X^{\frac{n}{n'}})$，上述求值等价于计算：

$\bm{(} \text{ }
M_Y((\omega^{J(0)})^{\frac{n}{n'}}), \text{ } M_Y((\omega^{J(1)})^{\frac{n}{n'}}), \cdots,  M_Y((\omega^{J(\frac{n}{2}-1)})^{\frac{n}{n'}}) \bm{)}$

$ $

$= \bm{(} \text{ }
M_Y((\omega^{\frac{n}{n'}})^{J(0)}), \text{ } M_Y((\omega^{\frac{n}{n'}})^{J(1)}), \cdots,  M_Y((\omega^{\frac{n}{n'}})^{J(\frac{n}{2}-1)})\bm{)}$

$ $

$= \bm{(} \text{ }
M_Y(\xi^{J(0)}), \text{ } M_Y(\xi^{J(1)}), \cdots,  M_Y(\xi^{J(\frac{n}{2}-1)})\bm{)}$

\textcolor{red}{ $\rhd$ 其中 $\xi = e^{\frac{i\pi}{n'}}$ 是本原 $(\mu=2n')$ 次单位根的基值与生成元}

$ $

注意 $\xi = \omega^{\frac{n}{n'}}$。因此，上述对 $M_Y(Y)$ 的求值，会将 $M_Y(Y)$ 在 $\frac{n'}{2}$ 个不同的本原 $(\mu=2n')$ 次单位根处的求值结果重复输出 $\dfrac{n}{n'}$ 次。



\subsection{模数自举}
\label{subsec:ckks-bootstrapping}

\textbf{- 参考文献：} \href{https://eprint.iacr.org/2018/153.pdf}{Bootstrapping for Approximate Homomorphic Encryption}~\cite{ckks}

在标准的分层求值安排中，重缩放最终会耗尽可用模数链。自举从较小模数 $q_0$ 开始，将其提升到工作模数 $q_L$，在内部消耗若干层后，返回模数为某个 $q_l>q_0$ 的密文。输出不一定处于最大模数 $q_L$。

取 $(A,B)\in R_{q_0}^2$ 的中心化代表元，并假设 $[B+AS]_{q_0}=\Delta M+E$。其未取模相位为
\[
B+AS=\Delta M+E+q_0K,
\]
其中 $K\in R$。这是关于整数提升的恒等式，并不是已经完成取模的解密恒等式。自举近似消除 $q_0K$ 项，保留原有数值误差，并引入新的误差。

\subsubsection{总体思路}
\label{subsubsec:ckks-bootstrapping-high-level}

自举的第一步，是直接将密文 $(A, B)$ 的模数从 $q_0$ 改为 $q_L$。此时，执行带模 $q_L$ 约减的解密，将输出：

$\textsf{RLWE}^{-1}\textbf{(}\textsf{ct} = (A, B)\textbf{)} \bmod q_L$

$ = B + A\cdot S \bmod q_L$

$= \Delta M + E + q_0 K \bmod q_L$

这种解释要求 $\|B+AS\|_\infty<q_L/2$。对于中心化输入代表元，有 $\|B+AS\|_\infty\leq(q_0/2)(1+\|S\|_1)$，因此可以通过选择足够大的 $q_L$ 保证该条件。仅指出秘密系数是三元值，并不能给出与维数无关的界；$\|S\|_1$ 可以达到 $n$。


在 $\Delta M + E + q_0 K \bmod q_L$ 中，$q_0K$ 项不再被模 $q_0$ 约减消去，因此密文解密所得明文多项式的每个 $i$ 次项都会得到被扰动的系数 $\Delta m_i + e_i + q_0\cdot k_i \bmod q_L$，而不是 $\Delta m_i + e_i \bmod q_L$。所以，现在需要消去各系数中的多余项 $q_0\cdot k_i \bmod q_L$，提取出纯粹的明文系数 $\Delta m_i + e_i$。



\begin{figure}[h!]
    \centering
  \includegraphics[width=1.0\linewidth]{figures/modulo-reduction-sine.png}
  \caption{正弦函数 $f(x) = \dfrac{q_0}{2\pi}\cdot \sin \left(\dfrac{2\pi x}{q_0}\right)$ 满足 $f(\Delta m_i + e_i + q_0k_i) \approx \Delta m_i + e_i$（假设 $\Delta m_i + e_i \ll q_0$）\href{https://eprint.iacr.org/2018/153.pdf}{（来源）}}
  \label{fig:modulo-reduction-sine}
\end{figure}

该函数以 $q_0$ 为周期。在 $x=kq_0$ 附近，它近似的是局部中心化余数 $x-kq_0$，而不是全局函数 $x$。它在 $kq_0$ 处的导数为 $1$。

$ $

结合这两个性质，对于输入 $x = \Delta m_i + e_i + q_0k_i$，有

\[
\begin{aligned}
f(\Delta m_i+e_i+q_0k_i)
&=\frac{q_0}{2\pi}\sin\left(\frac{2\pi(\Delta m_i+e_i+q_0k_i)}{q_0}\right)\\
&=\frac{q_0}{2\pi}\sin\left(\frac{2\pi(\Delta m_i+e_i)}{q_0}\right)\\
&\approx\Delta m_i+e_i.
\end{aligned}
\]

$ $

，前提是 $|\Delta m_i+e_i|/q_0$ 对于目标精度足够小。这是额外的自举参数条件，并非普通 CKKS 解密正确性的直接结果。令 $r_i=\Delta m_i+e_i$，正弦本身的近似误差满足
\[
\left|\frac{q_0}{2\pi}\sin(2\pi r_i/q_0)-r_i\right|
\leq\frac{2\pi^2}{3q_0^2}|r_i|^3.
\]
多项式近似和同态求值还会产生额外误差。由于密文乘法按槽位进行，算法先施加线性的系数到槽位变换，在槽位中求值非线性近似，再施加逆线性变换，流程见 \autoref{tab:ckks-bootstrapping-procedure}。

\begin{table}
\begin{tabular}{|lp{0.78\linewidth}|}
\hline
1 & \textsf{ModRaise}：将密文系数提升到 $q_L$，得到相位 $\Delta M+E+q_0K$。\\
2 & \textsf{CoeffToSlot}：同态施加逆嵌入，将 $n$ 个实相位系数放入两个密文。\\
3 & \textsf{EvalExp}：在明确的输入区间及误差预算下，求值周期模约减函数的非线性多项式近似。\\
4 & \textsf{SlotToCoeff}：将近似系数重新打包为一个密文，处于剩余模数 $q_l>q_0$，并包含全部累积误差。\\
\hline
\end{tabular}
\caption{CKKS 自举流程的高层描述}
\label{tab:ckks-bootstrapping-procedure}
\end{table}

\subsubsection{总体思路的数学展开}
\label{subsubsec:ckks-bootstrapping-high-level-correctness}

下述恒等式描述密文所表示的明文值，并不是求值过程中暴露秘密密钥或执行解密的操作。令 $T=\Delta M+E+q_0K$ 为提升后的相位，并仅将 $\mathsf{Enc}(T)$ 用作表示该相位的密文的简写。

\para{\textsf{\textbf{\underline{ModRaise}：}}} 将 $(A,B)$ 的系数提升到工作模数 $q_L$。在上述界成立时，提升后的密文具有原始相位 $T$。它可以直接作为同态求值的输入，无需在明文中计算 $AS$。

\para{\textsf{\textbf{\underline{CoeffToSlot}：}}} 令 $\sigma$ 为典范嵌入，$\pi$ 选出 $n/2$ 个独立槽位。实线性的逆嵌入将槽位 $\pi\sigma(T)$ 映为系数向量 $(T_i)_{i=0}^{n-1}$。其同态实现使用旋转、共轭、加法以及与编码矩阵常量的乘法。将 $n$ 个实系数存入两个密文，每个密文保存 $n/2$ 个系数，并记录这些变换的尺度与误差。

\para{\textsf{\textbf{\underline{EvalExp}：}}} 在所需输入区间上，同态求值近似 $f(x)=q_0\sin(2\pi x/q_0)/(2\pi)$ 的多项式电路 $\widehat f$。这是非线性计算，不是线性变换，也不是与单个编码多项式相乘。对于每个系数，
\[
\widehat f(T_i)=\Delta m_i+e_i+e_{b,i},
\]
其中 $e_{b,i}$ 包含正弦误差、多项式近似误差及全部求值误差。

\para{\textsf{\textbf{\underline{SlotToCoeff}：}}} 对所得两个密文施加逆线性打包操作。在所记录的输出尺度和模数 $q_l$ 下，最终相位表示为 $\Delta M+E+E_b$；$E_b$ 还包含最后一次变换的误差。

\para{\textsf{\textbf{\underline{无需秘密密钥的自举}：}}} 实际电路始终操作提升后的密文，不执行解密再加密；对于带噪声的随机化加密，也不存在 $\mathsf{Enc}(\mathsf{Dec}(\textsf{ct}))=\textsf{ct}$ 这样的恒等式。公开求值密钥支持旋转、共轭和重线性化。输出近似表示原始解码值，具有较大的可用模数，同时带有额外误差。

\subsubsection{细节：\textsf{CoeffToSlot}}
\label{subsubsec:ckks-bootstrapping-coefftoslot-details}

令 $\zeta_j=\omega^{J(j)}$，定义 $n/2\times n$ 求值矩阵 $U$，其中 $U_{j,k}=\zeta_j^k$。若 $\vec t$ 是 $T$ 的实系数向量，则其独立槽位为 $\vec z=U\vec t$。由单位根的正交性，
\[
\vec t=\frac{1}{n}(\overline U^{\,T}\vec z+U^T\overline{\vec z}).
\]
按系数索引的前半部与后半部，将 $U=(U_0\;U_1)$ 分为两个方阵块，则
\[
\vec t_0=\frac{1}{n}(\overline{U_0}^{\,T}\vec z+U_0^T\overline{\vec z}),\qquad
\vec t_1=\frac{1}{n}(\overline{U_1}^{\,T}\vec z+U_1^T\overline{\vec z}).
\]
同态施加这两个实线性映射，可得到槽位分别保存两半系数的密文，并产生编码和求值误差。共轭按 \autoref{subsec:ckks-conjugation} 计算；矩阵与向量乘法采用 \autoref{subsec:bfv-matrix-multiplication} 中的编码对角线和旋转方法。这些运算由密文与明文乘法及加法组成，层级消耗取决于选定的尺度安排。这里的记号不表示在明文中直接将矩阵乘以密文多项式的系数列表。

\subsubsection{细节：\textsf{EvalExp}}
\label{subsubsec:ckks-bootstrapping-evalexp-details}

利用正弦函数 $f(x) = \dfrac{q_0}{2\pi}\cdot \sin\left(\dfrac{2\pi x}{q_0}\right)$，通过计算 $f(\Delta m_i + e_i + q_0k_i) \approx \Delta m_i + e_i$，从 $\Delta m_i + e_i + q_0k_i$ 中近似消去 $q_0k_i$。当相对于 $q_0$ 而言，$\Delta m_i + e_i$ 非常接近 $x=0$（即 $\Delta m_i + e_i \ll q_0$）时，该近似有效。不过，对 $q_0k_i$ 的消去仍然是近似的（即 $\approx \Delta m_i + e_i$），因为在 $x=0$ 附近，$f(x)$ 为 $y \approx x$，而并非精确的 $y = x$。

这里的一个问题是，需要以 \textsf{ct\textsubscript{s1}} 与 \textsf{ct\textsubscript{s2}} 为输入（即 $f(\textsf{ct\textsubscript{s1}})$ 与 $f(\textsf{ct\textsubscript{s2}})$），同态求值 $f(x)$；但 FHE 仅支持 $(+, \cdot)$ 运算，而正弦曲线无法仅用 $(+, \cdot)$ 精确表示。因此，使用泰勒级数（\autoref{sec:taylor-series}）来近似正弦函数 $f(x)$：

$f(x) = f(a) + \dfrac{f'(a)}{1!}(x-a) + \dfrac{f''(a)}{2!}(x-a)^2 + \dfrac{f'''(a)}{3!}(x-a)^3 + \cdots = \sum\limits_{d=0}^{\infty}\dfrac{f^{(d)}(a)}{d!}(x-a)^d$

$ $

如果在 $x = 0$ 附近近似 $f(x)$，则近似多项式为：

$f(x) = \dfrac{q_0}{2\pi}\cdot \sin\left(\dfrac{2\pi}{q_0}\cdot 0\right) + \dfrac{q_0}{2\pi}\cdot\dfrac{2\pi}{q_0}\cdot\dfrac{\cos\left(\dfrac{2\pi}{q_0} \cdot 0\right)}{1!}\cdot x + \dfrac{q_0}{2\pi}\cdot \left(\dfrac{2\pi}{q_0}\right)^{2}\cdot\dfrac{-\sin\left(\dfrac{2\pi}{q_0} \cdot 0\right)}{2!}\cdot x^2 + \cdots$

$ = \dfrac{q_0}{2\pi}\cdot \sum\limits_{j=0}^{\infty}\left(\dfrac{(-1)^j}{(2j + 1)!}\cdot \left(\dfrac{2\pi x}{q_0}\right)^{2j + 1}\right)$

$ \approx \dfrac{q_0}{2\pi}\cdot \sum\limits_{j=0}^{h}\left(\dfrac{(-1)^j}{(2j + 1)!}\cdot \left(\dfrac{2\pi x}{q_0}\right)^{2j + 1}\right) = \hat f(x)$

$ $

其中，$\hat{f}(x)$ 是 $(2h + 1)$ 次多项式。

对于 $A,B$ 的中心化代表元及小秘密 $S$，整数相位满足 $B+AS=\Delta M+E+q_0K$。其系数范数至多为 $(q_0/2)(1+\|S\|_1)$，而对 $K$ 的界还包含 $\|\Delta M+E\|_\infty/q_0$。应选取 $k_{\max}$ 来约束整个归一化输入 $|x|/q_0$，包括残余消息和误差，而不只是 $|k_i|$。稀疏秘密可缩小该区间，但必须按其实际分布进行安全性分析。

归一化近似区间由 $\theta=2\pi x/q_0$ 决定，因此当 $|x|/q_0$ 固定时，单独增大 $q_0$ 不会扩大该区间。对于 $d$ 次多项式，霍纳求值使用 $d$ 次乘法，标准的并行幂构造可在多项式工作量下实现对数乘法深度。计算成本一般并不随次数指数增长。缩小输入后反复平方，是利用指数函数的特殊结构来减少工作量和深度：

\[
\begin{cases}
e^{i\cdot \theta} = \cos \theta + i \cdot \sin \theta\\
(e^{i \cdot \theta})^2 = e^{i\cdot 2\theta}\\
\end{cases}
\]

令 $\theta=2\pi x/q_0$。先在较小区间 $|\theta|/2^r\leq2\pi k_{\max}/2^r$ 上近似 $e^{i\theta/2^r}$，再将结果平方 $r$ 次，以近似 $e^{i\theta}$。所需泰勒次数取决于该归一化区间及目标最终误差。平方会放大初始近似误差并增加算术误差，因此初始容差必须考虑全部 $r$ 次平方。

在 $x = 0$ 附近使用 $d_0$ 次泰勒级数，可将 $e^{\frac{2\pi i x}{2^r q_0}}$ 近似为：

$f_e(x) = e^{\frac{2\pi i x}{2^r q_0}} \approx \sum\limits_{d=0}^{d_0}\dfrac{1}{d!}\left(\dfrac{2\pi i x}{2^rq_0}\right)^d = \hat{f}_e(x)$

$ $

然后，将 $\hat{f}_e$ 连续平方 $r$ 次，得到：

$(\hat{f}_e(x))^{2^r} \approx (f_e(x))^{2^r} = e^{i\frac{2\pi x}{q_0}} = e^{i\theta}$

$ $

接着，依据欧拉公式 $e^{i\cdot \theta} = \cos \theta + i \cdot \sin \theta$，可以推导如下关系：

$ $

$\overline{e^{i\cdot \theta}} = \cos \theta + \overline{i \cdot \sin \theta} $

$ e^{-i\cdot \theta} = \cos \theta - i \cdot \sin \theta$

$ e^{i\cdot \theta} - e^{-i\cdot \theta} = (\cos \theta + i \cdot \sin \theta) - (\cos \theta - i \cdot \sin \theta) = 2i\sin \theta$

$\sin \theta = \dfrac{-i}{2} \cdot (e^{i\cdot \theta} - e^{-i\cdot \theta})$

$\dfrac{q_0}{2\pi}\cdot \sin \theta = \dfrac{q_0}{2\pi}\cdot \dfrac{-i}{2} \cdot (e^{i\cdot \theta} - e^{-i\cdot \theta})$

$ $

代入 $\theta = \dfrac{2\pi x}{q_0}$，最终得到：

$\dfrac{q_0}{2\pi}\cdot \sin \left(\dfrac{2\pi x}{q_0}\right) = \dfrac{q_0}{2\pi}\cdot \dfrac{-i}{2} \cdot (e^{i\cdot \frac{2\pi x}{q_0}} - e^{-i\cdot \frac{2\pi x}{q_0}})$

$ $

利用上述最终关系，\textsf{EvalExp} 步骤按如下方式同态求值 $\dfrac{q_0}{2\pi}\cdot \sin \left(\dfrac{2\pi x}{q_0}\right)$ 的近似值，其中 $x = \Delta m_i + e_i + q_0k_i$：

\begin{enumerate}
\item 同态求值指数函数近似，得到 $\widehat z\approx e^{2\pi ix/q_0}$，并给出误差界。

\item 施加同态共轭（\autoref{subsec:ckks-conjugation}），得到 $\overline{\widehat z}$ 的近似，并计入密钥切换误差。

\item 将两者相减，再乘以编码常数 $-iq_0/(4\pi)$，以近似 $\frac{q_0}{2\pi}\sin(2\pi x/q_0)$。因子 $q_0/(2\pi)$ 不可省略；若仅乘以 $-i/2$，得到的将是未缩放的正弦。
\end{enumerate}

$ $

\textsf{EvalExp} 的结果是两个密文，其输入向量槽位存储了 $M_c$ 自举后的系数，即通过模约减 $q_0$ 将 $\Delta m_i + e_i + q_0k_i$ 转换为 $\Delta m_i + e_i + e_{bi}$。注意，$e_{bi}$ 是由以下三个因素产生的自举误差：（1）\textsf{CoeffToSlot}、\textsf{EvalExp} 和 \textsf{SlotToCoeff} 步骤固有的同态 $(+, \cdot)$ 运算噪声；（2）\textsf{EvalExp} 步骤中用泰勒多项式近似指数函数 $e^{i\theta}$ 产生的误差；（3）\textsf{EvalExp} 步骤的正弦曲线误差，因为曲线在 $x=0$ 附近并非精确等于 $y=x$，而只是 $y \approx x$。

注意，由于 \textsf{CoeffToSlot} 的输出被分为两个密文（\textsf{ct\textsubscript{s1}} 与 \textsf{ct\textsubscript{s2}}），\textsf{EvalExp} 的输出也包含两个密文：\textsf{ct\textsubscript{b1}} 与 \textsf{ct\textsubscript{b2}}。\textsf{ct\textsubscript{b1}} 的输入向量槽位存储 $(\Delta m_i + e_i + e_{bi})_{i=0}^{\frac{n}{2} - 1}$，而 \textsf{ct\textsubscript{b2}} 的输入向量槽位存储 $(\Delta m_i + e_i + e_{bi})_{i=\frac{n}{2}}^{n - 1}$。


\subsubsection{细节：\textsf{SlotToCoeff}}
\label{subsubsec:ckks-bootstrapping-slottocoeff-details}

采用 \autoref{subsubsec:ckks-bootstrapping-coefftoslot-details} 中相同的矩阵块 $U_0,U_1$，令两个 EvalExp 输出向量为 $\vec b_0,\vec b_1$，分别近似保存前半部与后半部系数。计算
\[
\vec z_{\mathrm{out}}=U_0\vec b_0+U_1\vec b_1.
\]
所得单个密文具有 $n/2$ 个复数槽位，表示的多项式具有由这两半拼接而成的 $n$ 个实系数。其隐含共轭槽位为 $\overline{\vec z_{\mathrm{out}}}$，并不是对同一系数多项式的另一次加密。同态变换会引入编码、密钥切换和尺度调整误差。因此，输出原始相位在模 $q_l$ 下为 $\Delta_{\mathrm{out}}M+E_{\mathrm{old,out}}+E_b$，其中全部尺度和误差都应沿具体实现传播。若计算安排恢复输入尺度，可以将其简写为 $\Delta M+E+E_b$。

\subsubsection{通过稀疏打包密文降低自举开销}
\label{subsubsec:ckks-bootstrapping-time-reduction}

假设 $2\leq n'<n$ 是一个真偶因子，令 $d=n/n'$。重复 $n'/2$ 个独立槽位会得到关于 $Y=X^d$ 的消息多项式，它有 $n'$ 个可能非零的系数。若起初只有一个非零槽位块、其余槽位均补零，可用 $\log_2 d$ 轮倍增来复制，每轮移位为 $-(n'/2)2^j$，其中 $0\leq j<\log_2d$。

模数提升后，仅有消息稀疏性还不够：相位 $\Delta M(Y)+E(X)+q_0K(X)$ 中的 $E$ 和 $K$ 一般并不稀疏。因此，原始 CKKS 自举论文第 5.2 节在 ModRaise 之后、CoeffToSlot 之前加入部分求和：
\[
\textsf{ct}\leftarrow\textsf{ct}+\textsf{Rotate}(\textsf{ct},(n'/2)2^j),
\qquad 0\leq j<\log_2d.
\]
这种类似迹映射的操作会消去次数不被 $d$ 整除的系数位置，并将保留系数乘以 $d$，同时产生求值误差。必须在记录的尺度和模约减周期中计入其归一化。若省略这一步，就不能据此丢弃提升后的其他系数。

较小的典范嵌入使用 $Y^{n'}+1$ 的根，具有 $n'$ 个实系数坐标，而不是 $n/n'$ 个。二次幂的真因子满足 $n'\leq n/2$，所以这些系数值可以装入一个密文。相应的较小系数与槽位变换以及模约减求值能够降低成本。实际收益取决于变换算法、旋转和归一化，并非所有参数下都恰好减半。

\subsubsection{总结}
\label{subsubsec:ckks-bootstrapping-summary}

\begin{tcolorbox}[title={\textbf{\tboxlabel{\ref*{subsubsec:ckks-bootstrapping-summary}} CKKS 自举}}]
\begin{enumerate}
\item 输入 $(A,B)\in R_{q_0}^2$，其中心化相位为 $\Delta M+E$，并满足所需输入区间界。
\item 将密文系数提升到 $q_L$。在模 $q_L$ 下满足不回绕界时，原始相位为 $\Delta M+E+q_0K$。RNS 实现需要计算相应的基扩展。
\item 执行 CoeffToSlot：完全打包输入使用两个密文，稀疏输入则采用上述正确归一化的流程。
\item 求值模约减的非线性多项式近似，同时预算近似误差与同态求值误差。
\item 执行 SlotToCoeff。输出处于剩余模数 $q_l>q_0$，并附有所记录的尺度和累积误差。采用等尺度简写时，其相位为 $\Delta M+E+E_b$，而不是精确的 $\Delta M+E$。
\end{enumerate}
\textbf{局限：} 自举不能恢复输入中已经损失的精度。后续数值计算仍须满足消息、误差和模数预算。
\end{tcolorbox}

\para{BFV 与 CKKS 自举的比较：} 精确方案的自举可通过求值精确解密函数消除先前的加密噪声，前提是满足其正确性条件。CKKS 会保留先前的数值近似误差，并增加新误差。无论固定模数的 BFV，还是一组固定 CKKS 参数，都不能脱离相应的噪声与精度条件支持任意无界工作。这里的 CKKS 流程直接使用提升后的密文；不需要单独的加密秘密内积，并不是因为所有中间尺度始终固定。

\subsubsection{降低自举噪声}
\label{subsubsec:ckks-bootstrapping-noise-reduction}

如 \autoref{subsubsec:ckks-bootstrapping-evalexp-details} 所述，自举过程会产生三类噪声：

%\renewcommand{\labelenumi}{\underline{Type-\arabic{enumi} Noise:}}
\begin{itemize}
\item \underline{第 1 类噪声}：\textsf{CoeffToSlot}、\textsf{EvalExp} 和 \textsf{SlotToCoeff} 步骤固有的同态 $(+, \cdot)$ 运算噪声
\item \underline{第 2 类噪声}：\textsf{EvalExp} 步骤中指数函数 $e^{i\theta}$ 的近似误差
\item \underline{第 3 类噪声}：\textsf{EvalExp} 步骤中的正弦曲线误差，即曲线在 $x=0$ 附近并非精确等于 $y=x$，而只是 $y \approx x$
\end{itemize}


$ $

对于选定的输入区间和参数，三类误差都必须有明确界。有限次多项式求值通常仍有非零近似误差，可以通过提高次数和计算精度来减小。改进模约减近似能减少正弦本身的误差，但不能省略对算术误差和近似误差的共同分析。

\begin{figure}[h!]
    \centering
  \includegraphics[width=1.0\linewidth]{figures/arc-sine.png}
  \caption{用于减小近似误差的反正弦曲线 \href{https://www.google.com/search?sca_esv=c744cb070de47b7e&sxsrf=ADLYWIKKaF93fcAFvJC4ehJso1E5qFTaow:1732546684173&q=arcsin(sin(x))&source=lnms&fbs=AEQNm0DmKhoYsBCHazhZSCWuALW8l8eUs1i3TeMYPF4tXSfZ95GzcfXnm5XYTvJV_9Qreh2py964ICpZJthXkELijctC8pFBYULoa3-fvQmwK0VJF0ntzsbN_W2CCJL9N57SWFNwWI58jCKaBJSdPgkprHQVK8H1PYOYWXMHTCCV-rDbC44rR6ANM870jZCZRtTKwFWtMIe2&sa=X&ved=2ahUKEwjZy6qt3_eJAxXbhlYBHXD2IvEQ0pQJegQIDxAB&cshid=1732546705369143&biw=1280&bih=635&dpr=1.5}{（来源）}}
  \label{fig:arc-sine}
\end{figure}

$ $

\para{$\bm{\arcsin(\sin(x))}$ 近似（\href{https://dl.acm.org/doi/10.1007/978-3-030-77870-5_22}{EUROCRYPT 2021}~\cite{cryptoeprint:2020/552}）：} 精确复合函数仅在 $[-\pi/2,\pi/2]$ 等主值区间上等于 $x$，其周期延拓的斜率交替变化。在所需的各个周期整数倍附近，它可以表示目标局部余数。同态实现需在指定区间上近似各组成函数，例如采用极小极大多项式。必须说明其近似误差、区间限制和额外深度；不能从该恒等式推出固定消耗三层的一般结论。

$ $

\para{Meta-BTS（\href{https://dl.acm.org/doi/10.1145/3548606.3560696}{CCS 2022}~\cite{cryptoeprint:2022/1167}）：} Meta-BTS 通过反复纠正前一次自举的误差来提高精度。其纠正思路可概括为：
\begin{enumerate}
\item 计算 $c_1=\mathsf{BTS}(c)$，表示 $m+\eta$。
\item 对齐模数与尺度，再构造 $d=\mathsf{ModDrop}(c_1)-c$，表示 $\eta$。应以自举后的密文减去原密文，才能使随后相减的符号正确。若不同时调整尺度，不能用重缩放替代 ModDrop。
\item 将小误差归一化到基础自举支持的输入区间，对其自举，再撤销归一化，得到表示 $\eta+\eta'$ 的 $\widehat d$。
\item 对齐参数，计算 $c_1-\widehat d$，表示 $m-\eta'$。
\end{enumerate}
归一化、精度及剩余层级是 Meta-BTS 定量保证的必要组成部分。上述概要并不声称，仅重复未经尺度调整的自举就一定能提高精度。

\para{噪声处理的局限：} 在满足算法的归一化与精度假设时，误差纠正可以减少自举新增误差。对固定的参数和实现，剩余算术误差与近似误差会限制可达到的精度。上述概要不能推出与输入无关的指数收敛或无限精度提升，也不能恢复原密文中已经损失的信息。


